\section{检测布局设计与实现}

\label{sec:detect}

\subsection{布局结构：原点 + 7 覆盖基点 + 12 均匀外圈点（20 点定案）}

由式~(\ref{eq:dircond})，检测区是半径 $R$ 的半圆盘 $H(g,\theta,R)$，$g,\theta,R$ 全部未知。
测量点集合须"尽量命中任意半圆盘"。按"减少路程 + 保证听率"双目标优化，最终定案为
\textbf{20 个测量位置}：

\begin{enumerate}
  \item \textbf{原点}：起点全频道扫描（位置成本为零），确认各频道在圆心处是否可听；
  \item \textbf{7 个覆盖基点}：直接复用问题三巡视站（\texttt{t3.config.SURVEY\_CENTERS}），
        全向覆盖保证（圆域内任意点到最近基点 $\le 1000$ m）；
  \item \textbf{12 个均匀方位外圈点}：半径 $1850$ m，方位间隔 $30^\circ$，补全定向源
        "迎光面"的命中。
\end{enumerate}

\begin{figure}[H]
  \centering
  \includegraphics[width=0.92\textwidth]{figures/fig_t4_sweep_lattice.pdf}
  \caption{检测扫描布局（20 点定案）：原点 + 7 覆盖基点 + 12 均匀外圈点；
  红色半圆盘为某贴边定向源的检测区示意（半径 $R=1000$ m、波束朝外），其内切圆内必有测量点
  命中。}
  \label{fig:lattice}
\end{figure}

\subsection{为什么外圈取 1850 m、12 个方位}

\paragraph{半径 1850 m。} 定向源的"迎光面"在源前方 $[|g|,\ |g|+R]$ 纵深；源最远半径
$1770$ m，故\textbf{任意贴边源}的迎光区内沿 $= 1770$ m。外圈设在其内侧 $1850$ m 时，对
任意半径 $R\in[1000,1500]$ 的源径向偏差 $< 80$ m，一次命中任意贴边源的迎光区。原
"7+21 中点外推"布局把多数点推到 $2270$ m，反而离内部半径源的迎光面太远（实测只剩
$99.81\%$）。

\paragraph{12 个方位（间隔 $30^\circ$）。} 相邻外圈点最大距离 $\approx 1850\cdot\sin
15^\circ \approx 479$ m $\ll$ 接收半径 $1000$ m 的一半，任意方向的源都能被方向差
$\le 15^\circ$ 的外圈点听到。曾试 $11$ 点（间隔 $32.7^\circ$、布局路线短 $294$ m），但方位
间隔变宽使\textbf{顺路清除的方位扇区也变宽}、绕路更多，20 局实测总里程反升 $464$ m——故
$12$ 点 $\times 1850$ m 是"减少路程"的实战极值。

\subsection{为什么不加内部补点（20 点定案）}

曾试过在基点带之内加 $3$ 个内部补点（半径 $850$ m、间隔 $120^\circ$、旋转 $15^\circ$，
23 位置）：把内部方向空隙从 $\sim 51^\circ$ 收窄到 $\sim 26^\circ$（文献"中心 + 多环带"
结构），听率从 $99.9891\%$ 提到 $99.9974\%$。但\textbf{20 局演练实测反而慢 $243$ s}（测向
多约 $38$ 次/局，$6445 \to 6688$ s），而 20 点版 20 局实测 $256/256$ 全清、完成率无差异。
经第~\ref{sec:verify} 节神经网络耗时寻优交叉验证（找不到更省时布局）后，\textbf{最终定案
20 点版}：耗时与任务完成率双优。

\subsection{听到率统计（422.5 万例全口径）}

调用 \texttt{cumcm.t4.sweep.verify\_hearing\_stats}，对 20 点定案做三类合计
$4\,224\,640$ 例校验（表~\ref{tab:hear}）：真实随机场景蒙特卡洛 $400$ 万例、贴边对抗
（源贴 $1770$ m、波束径向 $\pm 8^\circ\times 360$ 方位 $\times 3$ 档 $R$）$43\,200$ 例、
精细对抗（$g$ 半径 $0\sim1770 \times 48$ 方位 $\times 180$ 波束 $\times 3$ 档 $R$）
$181\,440$ 例。

\begin{table}[H]
  \centering
  \caption{20 点定案布局的听到率统计：合计 $4\,224\,640$ 例漏 $460$（听到率
  $99.9891\%$），比"7+21 中点外推"（$99.81\%$）高两个数量级且贴边对抗 0 漏。}
  \label{tab:hear}
  \begin{tabular}{lccr}
    \toprule
    \textbf{校验} & \textbf{算例} & \textbf{漏} & \textbf{听到率} \\
    \midrule
    蒙特卡洛（真实随机场景） & $4\,000\,000$ & $320$ & $99.9920\%$ \\
    贴边对抗（最严，外翻光束刀口） & $43\,200$ & $0$ & $100.000\%$ \\
    精细对抗（内部半径源全枚举） & $181\,440$ & $140$ & $99.9228\%$ \\
    \textbf{合计} & \textbf{$4\,224\,640$} & \textbf{$460$} & \textbf{$99.9891\%$} \\
    \bottomrule
  \end{tabular}
\end{table}

残余漏例集中在"内部半径源 + 波束朝外 + 方向恰落 7 基点方位空隙"（只有内层测量点可命中），
非严格保证，论文按实测统计报告（数据来源 \texttt{results/t4/t4\_sweep\_plan.json} 的
\texttt{verification} 字段）。

\subsection{访问路线}

以原点为起点的开放路径，由最近邻 + 2-opt 局部搜索精修（\texttt{common.routing}，确定性
无随机算子），里程 $16\,996$ m（比"7+21 中点外推"的 $21\,580$ m 省 $21\%$，还少 $9$ 个
位置）。顺序与里程不影响检测效果（只取决于点集），只影响行驶耗时。